%% ===========================================================================
%% MASTER CV - Marcus D. Yang
%% ---------------------------------------------------------------------------
%% This file is the FROZEN FORMATTING STANDARD and the FACT SOURCE for every
%% tailored CV. Two rules:
%%   1. The preamble below (geometry, font, \section rule, \erow macro, list
%%      spacing) is frozen. A tailored CV COPIES it verbatim - never retune it.
%%   2. Everything between \begin{document} and \end{document} is CONTENT and
%%      is fully re-orderable / re-groupable per the job description. See
%%      .claude/skills/job-application-assistant/05-cv-templates.md.
%% Compile: lualatex -interaction=nonstopmode main_example.tex   (MUST be 1 page)
%% Visual reference: cv/Marcus Yang Master Resume.pdf
%% ===========================================================================
\documentclass[11pt,letterpaper]{article}

% --- Fonts: Times New Roman clone, bundled with every modern TeX dist -------
\usepackage{fontspec}
\setmainfont{TeX Gyre Termes}
\usepackage{microtype}
\linespread{0.93}   % match the master's tight (Word "single") leading

% --- Page geometry (master: 0.8in sides, ~0.5in top) ----------------------
\usepackage[top=0.5in,bottom=0.5in,left=0.8in,right=0.8in]{geometry}

% --- Section headings: bold, full-width rule underneath --------------------
\usepackage{titlesec}
\titleformat{\section}{\fontsize{12.5}{14}\selectfont\bfseries}{}{0pt}{}[{\vspace{1pt}\hrule height 0.6pt}]
\titlespacing*{\section}{0pt}{6pt}{2pt}

% --- Tight bullet lists ---------------------------------------------------
\usepackage{enumitem}
\setlist[itemize]{leftmargin=1.35em,itemsep=1pt,topsep=1pt,parsep=0pt,
                  partopsep=0pt,label=\textbullet}

\usepackage[hidelinks]{hyperref}
\pagestyle{empty}
\setlength{\parindent}{0pt}

% --- \erow{left}{right}: entry/degree row, left flush-left, right flush-right
%     Caller supplies ALL markup in {left} (\textbf{...}, \textit{...}, etc.).
\newlength{\erowdate}\setlength{\erowdate}{2.0in}
\newcommand{\erow}[2]{%
  \noindent\makebox[\dimexpr\linewidth-\erowdate\relax][l]{#1}%
  \makebox[\erowdate][r]{#2}\par}

% --- \entrygap: vertical space between entries within a section -----------
\newcommand{\entrygap}{\vspace{2pt}}

\begin{document}

% ===========================================================================
%  HEADER
% ===========================================================================
\begin{center}
  {\fontsize{17}{20}\selectfont\bfseries Marcus D. Yang}\\[2pt]
  {\footnotesize (650) 954-0288 \textbar{} \href{mailto:yangdmarcus@gmail.com}{yangdmarcus@gmail.com} \textbar{} \href{https://www.linkedin.com/in/marcusdyang}{www.linkedin.com/in/marcusdyang} \textbar{} \href{https://projects-website-psi.vercel.app/}{https://projects-website-psi.vercel.app/}}
\end{center}

% ===========================================================================
%  EDUCATION
% ===========================================================================
\section{Education}
\erow{\textbf{University of Illinois Urbana-Champaign} \textbar{} Champaign, IL}{August 2024 - May 2028}
\erow{\textit{B.S. in Aerospace Engineering with a Computer Science minor, James Scholar Honors}}{GPA: 3.57/4}
\textbf{Relevant Coursework:} Aerospace Control Systems, Electrical and Electronic Circuits, Linear Algebra, Discrete Structures, Intro to Computer Science II, Dynamics, Aerospace Structures, Thermodynamics

% ===========================================================================
%  ORGANIZATIONAL AND TECHNICAL EXPERIENCE
% ===========================================================================
\section{Organizational and Technical Experience}

\erow{\textbf{TVC Rocket} (Python, C++, Embedded Systems)}{Summer 2026}
\begin{itemize}
  \item Integrated a Teensy 4.1 flight computer, IMU, and servo actuators with a 2-axis TVC gimbal, avionics bay, and recovery system on a commercial rocket body.
  \item Engineered a closed-loop PID flight-control system in C++ on a Teensy 4.1 microcontroller, fusing IMU data to steer the gimbal in real time.
  \item Built a Python 6-DOF flight simulation and ran Monte Carlo sensitivity sweeps of sensor noise, motor burn variance, and wind across hundreds of trials to validate the control algorithm.
  \item Found from the sweep data that landing performance was far more sensitive to IMU accelerometer bias than servo speed, redirecting effort toward sensor calibration.
  \item Ran hardware-in-the-loop integration testing of the PID firmware against a live IMU feed to verify gimbal servo response before static fire.
  \item Root-caused two static fire test failures by analyzing telemetry logs and frame-by-frame video, isolating a misreferenced attitude axis and an inverted control loop sign.
  \item Logged onboard telemetry at 83 Hz with zero dropouts across two test campaigns (15,728 and 12,611 samples).
\end{itemize}
\entrygap

\erow{\textbf{Liquid Rocketry at Illinois}}{Spring 2025 - Present}
\begin{itemize}
  \item Modeled P\&ID system architecture in Siemens NX and integrated components into the pressurant tank assembly, sourcing and constraining CAD hardware to spatial requirements.
  \item Assessed pressurant tank options, produced mechanical integration models (CAD) and verified fit against vehicle structural constraints and P\&ID.
  \item Engineered a KiCad PCB radio transmitter (USB-C + MCU) prototype enabling ground-test telemetry links.
  \item Created CAD geometries for canard-grid fin configurations and ran ANSYS Fluent CFD cases to quantify aerodynamic interactions across 0--20° deflection.
\end{itemize}
\entrygap

\erow{\textbf{Model Rocket Project} (Siemens NX, OpenRocket)}{Fall 2024}
\begin{itemize}
  \item Designed custom fin geometry in Siemens NX and fabricated a physical prototype to hit a target apogee.
  \item Validated performance against OpenRocket simulations; flight reached 430 ft apogee with 25 ft downrange drift.
\end{itemize}

% ===========================================================================
%  TECHNICAL SKILLS
% ===========================================================================
\section{Technical Skills}
\textbf{Computer Languages:} Python, C++, Java, JavaScript\\
\textbf{Software:} Siemens NX, Fusion 360, Eagle PCB, KiCad, OpenRocket, XFLR5, Ansys Fluent, Excel\\
\textbf{Shop Tools:} Bambu A1 Mini, Spot Welder, Belt Sander, Laser Cutter, CNC Plasma Cutter, Table Saw, Drill Press\\
\textbf{Materials:} Epoxy Resin, Fiberglass, Sheet Metal, Aluminum Pipes, Plywood, PETG, PLA

% ===========================================================================
%  LEADERSHIP AND ACTIVITIES
% ===========================================================================
\section{Leadership and Activities}

\erow{\textbf{Asian American Association}}{Fall 2024 - Present}
Treasurer \textbar{} Fundraising Officer
\begin{itemize}
  \item Managed membership operations for 300+ members, coordinating records and organizational logistics.
  \item Secured partnerships with 15 on-campus sponsors to support events and fundraising initiatives.
\end{itemize}
\entrygap

\erow{\textbf{Flight Club Aerospace}}{Summer 2021 - Spring 2024}
Landing Gear Research Lead \textbar{} CAD \textbar{} Fabrication
\begin{itemize}
  \item Manufactured wing and control surfaces: cut foam ribs, trimmed aluminum spars with a circular saw, and applied epoxy \& fiberglass layups to shape aerodynamic skins; aligned two 7-ft wings for flight testing.
  \item Drafted gussets and ribs in Fusion 360 and confirmed fit-checks between structural and aerodynamic components.
  \item Oversaw logistics for wing transport/disassembly/reassembly and assembled and presented landing-gear recommendations based on structural load and cost analyses.
\end{itemize}

\end{document}
